\documentclass[11pt]{article}
\usepackage[margin=1in]{geometry}
\usepackage{amsmath,amssymb,amsthm,enumitem,hyperref}
\hypersetup{colorlinks=true,urlcolor=blue,linkcolor=blue}
\newtheorem{theorem}{Theorem}
\title{A rational counterexample to pentagon barycenter coincidence}
\author{Principia Math In Progress: MathDB 360860}
\date{21 September 2026}
\begin{document}
\maketitle
\begin{abstract}
We give a strictly convex pentagon with rational vertices on the unit circle
whose hyperbolic barycenter and Euclidean least-squares point coincide away
from the origin, although its five consecutive cross-ratios are not equal.
This disproves Conjecture 36 of Arnold and Arreche.
\end{abstract}

\section*{Task 1: the original problem}
We use the notation of the actual source \cite{AA}: $S_{\mathbf P}$ is the
hyperbolic barycenter and $L_{\mathbf P}$ minimizes the sum of squared
Euclidean distances to the side lines. Both are expressed in the
Beltrami--Cayley--Klein unit-disc coordinates. Conjecture 7.8 of the preprint,
retained as Conjecture 36 of the journal final, asserts that
$S_{\mathbf P}=L_{\mathbf P}$ for a pentagon if and only if its five
consecutive cross-ratios are equal or $S_{\mathbf P}=(0,0)$.

\begin{theorem}
Let the vertices of $\mathbf P$, in cyclic order, be
\[
 P_1=(1,0),\quad P_2=\left(\frac7{25},\frac{24}{25}\right),\quad
 P_3=\left(-\frac45,\frac35\right),\quad
 P_4=\left(-\frac45,-\frac35\right),\quad
 P_5=\left(\frac7{25},-\frac{24}{25}\right).
\]
Then $\mathbf P$ is strictly convex and
\[
                     S_{\mathbf P}=L_{\mathbf P}
                        =\left(-\frac1{62},0\right),
\]
but its five consecutive cross-ratios are not all equal.
\end{theorem}

\section*{The polygon and the source barycenter formula}
The source's stereographic parameters give
\[
 P(p)=\left(\frac{1-p^2}{1+p^2},\frac{2p}{1+p^2}\right),\qquad
 (p_1,p_2,p_3,p_4,p_5)=\left(0,\frac34,3,-3,-\frac34\right).
\]
These are five distinct points on the unit circle in counterclockwise order,
so their polygon is strictly convex. With indices modulo five, the source
defines
\[
 I=\sum_i\frac1{p_i-p_{i+1}},\qquad
 J=\frac12\sum_i\frac{p_i+p_{i+1}}{p_i-p_{i+1}},\qquad
 K=\sum_i\frac{p_ip_{i+1}}{p_i-p_{i+1}}.
\]
Its equation (3.3) and Theorem 5 in the journal final give
\[
 S_{\mathbf P}=\left(\frac{I-K}{I+K},\frac{2J}{I+K}\right).
\]
Direct substitution gives $I=-61/18$, $J=0$, $K=-7/2$, and hence
$I+K=-62/9$ and $I-K=1/9$. Therefore $S_{\mathbf P}=(-1/62,0)$.

\section*{The Euclidean least-squares point}
The five side lines, in order, are
\[
 4x+3y=4,\quad -5x+15y=13,\quad 5x=-4,\quad
 -5x-15y=13,\quad 4x-3y=4.
\]
Thus their squared-distance sum is
\begin{align*}
 F(x,y)&=\frac{(4x+3y-4)^2+(4x-3y-4)^2}{25}
          +\frac{(5x+4)^2}{25}\\
 &\quad+\frac{(-5x+15y-13)^2+(-5x-15y-13)^2}{250}\\
 &=\frac{62}{25}x^2+\frac{63}{25}y^2+\frac2{25}x+\frac{409}{125}.
\end{align*}
Its Hessian is $\operatorname{diag}(124/25,126/25)$, which is positive
definite. Its unique global minimum is consequently at $(-1/62,0)$.
This proves $L_{\mathbf P}=S_{\mathbf P}\ne(0,0)$.

\section*{The unequal cross-ratios}
To test equality, fix the ordered convention
\[
 R(a,b,c,d)=\frac{(a-c)(b-d)}{(a-d)(b-c)},\qquad
 R_i=R(p_i,p_{i+1},p_{i+2},p_{i+3}).
\]
The five values are
\[
 (R_1,R_2,R_3,R_4,R_5)
       =\left(\frac53,\frac{25}{16},\frac53,\frac85,\frac85\right).
\]
In particular $R_1\ne R_2$. Changing a fixed ordered cross-ratio convention
applies an injective fractional-linear function (possibly with a change
of sign) to each value, so it does not affect whether the five values
are equal. This disproves the necessary direction of the conjecture.
\hfill$\square$

\section*{An independent hyperbolic check}
For a chord line $ax+by=c$ in the unit-disc model, the hyperbolic distance
$d_H$ satisfies
\[
 \sinh d_H((x,y),\{ax+by=c\})
 =\frac{|ax+by-c|}{\sqrt{a^2+b^2-c^2}\sqrt{1-x^2-y^2}}.
\]
This follows directly from the polar and hyperboloid formulas (2.5)--(2.7)
of \cite{AA}. Choose the inward linear forms
\[
 q_1=4-4x-3y,\quad q_2=13+5x-15y,\quad q_3=4+5x,\quad
 q_4=13+5x+15y,\quad q_5=4-4x+3y.
\]
The corresponding normalizing factors are $(3,9,3,9,3)$, so throughout
the open disc the hyperbolic objective $G=\sum_i\sinh d_{H,i}$ obeys
\[
 G(x,y)\ge \frac{62+x}{9\sqrt{1-x^2-y^2}}
             \ge\frac{\sqrt{3843}}9.
\]
The second inequality, including its equality case, follows from
\[
 (62+x)^2-3843(1-x^2-y^2)=(62x+1)^2+3843y^2
\]
and $62+x>0$. At $(-1/62,0)$ all five $q_i$ are positive and both
inequalities are equalities. This independently certifies the same unique
hyperbolic minimum, with no appeal to a numerical optimizer.

\section*{Scope and verification}
The example satisfies convexity, cyclic order, distinctness, and the unit
circumcircle normalization. It refutes the source statement itself; the
abbreviated MathDB names for the two centers are not used as a loophole.
The companion verifier uses exact arithmetic for the geometry, both center
calculations, the hyperbolic minimum certificate, and the cross-ratios.
Review is by the authoring agent only, without independent peer review or
a proof-assistant certificate.

\section*{Proposed follow-up Tasks 2--5}
\begin{enumerate}[label=Task \arabic*:,start=2,leftmargin=4.6em]
\item Classify coincidence in the reflection-symmetric family
$p=(0,a,b,-b,-a)$ with $0<a<b$.
\item Determine the connected components of all strictly convex pentagons
with $S_{\mathbf P}=L_{\mathbf P}\ne0$, modulo rotations of the unit disc.
\item Characterize which strictly convex ideal pentagons with
$S_{\mathbf P}=0$ also have $L_{\mathbf P}=0$.
\item For every $n\ge6$, determine whether a strictly convex rational
unit-circle $n$-gon exists with coincident nonzero centers and unequal
consecutive cross-ratios.
\end{enumerate}

\begin{thebibliography}{1}
\bibitem{AA} M. Arnold and C. E. Arreche,
\emph{Symmedians as Hyperbolic Barycenters},
Comptes Rendus Math\'ematique \textbf{362} (2024), 1743--1762.
\href{https://doi.org/10.5802/crmath.677}{Journal version (DOI)}.
Preprint: \href{https://arxiv.org/abs/2311.14194}{arXiv:2311.14194v1}.
\end{thebibliography}
\end{document}
